\section{Discussion and Limitations}
\label{sec:discussion}

\paragraph{Random vectors do not yield a significant gain.}
A central concern for any per-session personalization mechanism is that the bilinear term $h^\top A u$ may merely memorize session identity rather than model user preference: with a random per-session $u$, the 12{,}288 parameters of $A$ could in principle fit session-specific offsets. Our random-vector control (Table~\ref{tab:representation}) directly tests this. At $n{=}201$ test sessions, the random variant is \emph{not} significant ($\Delta{=}{+}0.027$, CI $[-0.010, +0.065]$), while behavioral features ($+0.084$) and LLM schemas ($+0.054$) are. This argues against pure session-identity memorization: if the bilinear term only needed a per-session code, random vectors should have worked equally well. The gain requires genuine user signal.

\paragraph{Behavioral features beat LLM schemas.}
We initially hypothesized that LLM-extracted typed schemas, being semantically richer, would outperform six-dimensional behavioral session features. They do not: behavioral features yield a larger and tighter gain ($+0.084$ vs $+0.054$). One explanation is that the LLM schema (extracted by a Phi-3-mini model under a closed-vocabulary constraint) introduces noise that the simpler, directly-computed behavioral features avoid. Another is that for the evidence-reranking task on AOL, behavioral statistics (query overlap, history length, reformulation) carry more signal than preference-type labels do. The LLM schema retains value for interpretability --- it produces a human-readable four-dimensional preference label --- but it is not the strongest user representation for ranking quality. We treat this as an honest negative finding on the LLM-extraction hypothesis.

\paragraph{Sample size.}
The test set of $n{=}201$ sessions provides adequate statistical power to detect the behavioral-features gain ($+0.084$, CI $[+0.034, +0.131]$) and the LLM-schema gain ($+0.054$, CI $[+0.013, +0.094]$), both with CIs excluding zero. The random-vector non-significance ($+0.027$, CI $[-0.010, +0.065]$) is also informative: it rules out the session-identity confound at this sample size. The difference between behavioral ($+0.084$) and LLM-schema ($+0.054$) gains has overlapping CIs, so we do not claim behavioral features are \emph{proven} superior --- only that they are at least as good and more parsimonious.

\paragraph{Data split protocol and leakage safeguards.}
Sessions are split into train/dev/test by anonymized user id (70/15/15), ensuring no user appears in more than one split. Within each split, training pairs are constructed from position-debiased clicks (rank $> 3$), and negative samples are drawn from cross-session clicks within the same split, excluding any URL in the current user's history. This protocol prevents user leakage, topic leakage, and direct negative reuse across splits.

\paragraph{AOL click-only supervision.}
AOL-IA provides click metadata but no full-SERP impressions; our negatives are cross-session clicked distractors rather than documents a user saw and skipped. This is a standard proxy in click-log IR, but it means absolute nDCG values are not comparable to doc-aware benchmarks. The \emph{delta} from personalization is the contribution, not the absolute score.

\paragraph{Backbone choice.}
The BM25+BGE query-flow backbone is weaker than a doc-aware cross-encoder such as monoT5. We chose it because AOL-IA does not ship document text, and downloading the Internet-Archive snapshot of 1.5M pages was infeasible within our timeline. The bilinear intervention is backbone-agnostic and transfers directly to any frozen scorer that exposes $s_{\mathrm{base}}$ and $h$; we expect absolute gains to differ on a doc-aware backbone.

\paragraph{Seed sensitivity.}
We report the \texttt{classic} variant under three independent seeds (Table~\ref{tab:seed}): $\Delta \in \{+0.072, +0.084, +0.087\}$, mean $+0.081$, std $0.008$, with every seed's 95\% CI excluding zero. The gain is stable to initialization and is not an artifact of $\text{seed}{=}42$.

\paragraph{Candidate-set construction caveat.}
The candidate set in our evaluation is one positive plus nine cross-session clicked distractors for the same query. This is a controlled benchmark, not a faithful end-to-end deployment distribution: a retrieval stage over a larger document pool, with non-clicked URLs drawn from the same user's SERP, would be a more realistic test, and a sensitivity analysis varying distractor construction (same-query vs.\ BM25 hard negatives, varying pool sizes) is a prerequisite for stronger causal claims. We treat the reported gains as evidence that a useful personalization signal exists in this benchmark, not as a definitive estimate of deployment performance.

\paragraph{Mechanism comparison.}
Section~\ref{sec:exp:mech} compares the bilinear form against concatenation-plus-linear and FiLM under matched user representation. The bilinear form is Pareto-optimal: it beats concatenation ($+0.084$ vs $+0.054$) and matches FiLM ($+0.081$) at roughly half the parameters ($12{,}304$ vs $26{,}112$). The bilinear form is not uniquely necessary --- FiLM works too --- but it is the most parameter-efficient way to reach the plateau on this benchmark. The user-vector dimension saturates at $u_{\dim}{=}8$ (Sec.~\ref{sec:exp:dim}), so $u_{\dim}{=}16$ is a safe default rather than a tuned value.

\paragraph{Interpretability vs. ranking quality.}
The LLM-extracted typed schema does not improve nDCG over behavioral or random vectors, but it is the only construction that yields a human-readable preference explanation. For deployments that require ``why was this evidence shown to this user,'' the schema retains independent value. We treat ranking quality and interpretability as orthogonal axes and do not claim the schema dominates on both.

\paragraph{Threats to validity.}
Three threats stand out. (i) Single LLM (GPT-5.4-mini) for schema extraction; a different model could change the schema-quality parity. (ii) The bilinear term's 12{,}288 parameters may overfit session patterns at the $n=15$ test scale; we mitigated this with dev-based early stopping but cannot fully rule it out. (iii) User-disjoint splitting prevents user leakage but does not prevent topical leakage across sessions.
